\section{Ondes de surfaces circulaires}

\subsection{Description des manipulations}

La deuxième expérience se scinde en deux parties. La première permet de vérifier les bonnes conditions de mesures de l'installation afin qu'en un second temps permettre d'étudier les interférences.

\begin{figure}[h!]
    \centering
    % Première figure
    \begin{minipage}[b]{0.6\linewidth}
    \centering
    \includegraphics[width=\linewidth]{manip_cuve.png}
    \caption{$Z_L = 2000 \ohm$}
    \label{fig:r2000}
    \end{minipage}
    \hfill
    % Deuxième figure
    \begin{minipage}[b]{0.3\linewidth}
         L'installation visible sur Fig. \ref{fig:montage_cuve} contient donc d'abord un exitateur (dispositif pneumatique avec transmissions des oscillations accoustiques par un conduit souple) pour mesurer les longueurs d'ondes de surface. Un deuxième est nécessaire si l'on souhaite mesurer des interférences (entre deux sources par définition). Le \\\\

    \centering
    \end{minipage}
\end{figure}
Nous commençons par utiliser un niveau à bulle pour nous assurer que la vitre de la cuve soit la plus horizontale possible et éviter toute distorsion des ondes. Une fois la cuve correctement nivelée, nous la remplissons d'une solution d'eau savonneuse de près de 1cm. Nous placons ensuite un excitateur d’ondes sur son support, ajustons la hauteur de l’excitateur afin qu'il soit partiellement immergé et transmette par l'extrémité de la buse les oscillations générées. Cela assure que les ondes soit proprement propagées sur la surface. Enfin, nous choisissons une dizaine de fréquences d’excitation (basses fréquences) comprises entre 15 et 80 Hz que nous lisons sur un oscilloscope numérique. Les impulsions entre stroboscope et l'exitateur sont synchronisés au travers les différentes fréquences. Cela permet de "geler" les ondes et de mesurer précisément les longueurs d’ondes sur un écran translucide en fonction de la fréquence générée. 
\\\\
En un second temps, un deuxième exitateur est placé à une distance fixe en veillant à ce que la longueur de tube soit identique dans les deux buses. Cela permet d'éviter un déphasage entres les deux sources. Nous mesurons ensuite au travers de l'écran les interférences générées en fonction de plusieurs basses fréquences également autour de 15 et 80Hz. Nous prenons pour chaque fréquence une photo à ~2.5m de l'écran pour la traiter l'image et déterminer les hyperboles correspondantes.
